\documentclass[a4paper]{article}
\usepackage[margin=0.75in]{geometry}
\usepackage{parskip}
\usepackage{amsmath}
\usepackage{amsxtra}
\usepackage{amssymb}
\usepackage{mathtools}
\usepackage{multicol}
\usepackage{float}
\usepackage{multirow}
\usepackage{ragged2e}
\usepackage{booktabs}
\usepackage{tabularx}
\usepackage{graphicx}
\usepackage{subcaption}
\usepackage[colorlinks=false, filecolor=red]{hyperref}

\newcommand{\since}{\because}
\newcommand{\R}{\mathbb{R}}
% \mathtoolsset{showonlyrefs}
\numberwithin{equation}{section}
\numberwithin{figure}{section}
\numberwithin{table}{section}

\title{Bead on a Rotating Hoop: Analysis of Chaos}
\author{Roshan Mathew Philip and Ashlin V. Thomas}
\date{\today}
\begin{document}
\maketitle
\raggedright

\section{Bead on a Hoop}
The lagrangian of our system is:
\begin{equation}
    L = \frac{1}{2}mR^2(\dot{\theta}^2 + \omega^2 \sin^2\theta) + mgR(\sin\theta \sin\alpha \cos\omega t + \cos\alpha \cos\theta)
\end{equation}
We get the equations of motion:
\begin{align}
    \ddot{\theta} &= -\gamma \dot{\theta} - \left(g_{eff} \cos\alpha - \omega^2 \cos\theta\right) + g_{eff}\sin\alpha \cos\omega t \cos\theta\\
    \intertext{Where $g_{eff} = g/R$. Using small $\alpha$ approximation}
    \ddot{\theta} & +\gamma \dot{\theta} + (g_{eff} - \omega^2 \cos\theta)\sin\theta = g_{eff}\alpha\cos\omega t\cos\theta
    \intertext{Using small $\theta$ approximation}
    \ddot{\theta} & +\gamma \dot{\theta} + (g_{eff} - \omega^2 )\theta = g_{eff}\alpha\cos\omega t
    \label{eq:small_ang}
\end{align}

We see that, at small angles, the system behaves like a damped driven oscillator;
with a natural frequency $\sqrt{g_{eff} - \omega^2}$. Note that the natural
frequency depends on the driving frequency.
Since it is a dampled driven oscilllator we get a resonance
at
\begin{align}
    (g_{eff}-\omega_{res}^2) - \omega_{res}^2 &= 0\\
    \implies \omega_{res} &= \sqrt{\frac{g_{eff}}{2}}
\end{align}

We use eq$\eqref{eq:small_ang}$ near resonance only to show the existence of
resonance. We can't use it for evaluating the trajectories or even
the amplitude of the resonant oscillations since $\theta$ would be large.

See \emph{Lisandro A Raviola et al 2017 Eur. J. Phys. 38} for 
the correction to use when we consider a bead rotating without slipping
instead of a point particle.

\section{Hamiltonian Chaos}
Because of Liouville's theorem chaos is hard to achieve in Hamiltonian systems.
We have two major explanations for why chaos emerges in Hamiltonian systems,
KAM theory, and Melnikov theory.

\subsection{Action-Angle Variables}
Using the Hamilton-Jacobi theory we get our generalised coordinates
in a very particular form $(I_\theta,\theta)$ where $\theta$ is an angle and
$I_\theta$ is the corresponding momentum. Furthermore, the generalised
coordinate $\theta$ is cyclic ($\therefore I_\theta$ is contant), and $\dot{\theta}=\omega(I_\theta)$.\\

Therefore, in the phase space $I_\theta$ is fixed while $\theta$ vary.
The full phase space is isomorphic to $\R^n\times T^n$ where $n$ is the
number of generalised angles, and $T^n$ is an n-Torus:
$$T^n = S^1 \times \dots (\text{$n$ times})\dots \times S^1$$
\subsection{Kolmogorov–Arnold–Moser (KAM) theory}
We start with an integrable Hamiltonian system ($H_0$) and a perturbation
to it that makes it non-integrable.\\
Note $H_0\equiv H_0(I+i)$
\begin{align}
    \dot{I_i}= 0; \quad&\quad \dot{\theta_i}= \omega_i(I)\\
    \omega_i(I)=&\frac{\partial H_0}{\partial I_i} 
    \intertext{Now perturb the Hamiltonian}
    H(I, \theta) &= H_0(I) + \epsilon H_1(I,\theta)
\end{align}

The question KAM theory tries to answer is whether the invariant torii survive.

\subsubsection{Resonances}
A resonance occurs when
$$k\cdot\omega = 0$$
for some vector $k=0$ with integer components. Or in other words, the
ratios between the frequencies of the different $\theta$ are rational.

\end{document}
